\documentclass[11pt,a4paper]{amsart}

\usepackage[T1]{fontenc}
\usepackage{lmodern}
\usepackage[protrusion=true,expansion=false]{microtype}
\usepackage[margin=30mm]{geometry}
\usepackage{amsmath,amssymb,amsthm,mathtools}
\usepackage[dvipsnames]{xcolor}
\usepackage[colorlinks=true,linkcolor=MidnightBlue,citecolor=MidnightBlue,urlcolor=MidnightBlue]{hyperref}
\hypersetup{
 pdftitle={A Boundary-Constant Obstruction to Real-Rooted Riemann Incomplete-Gamma Approximants},
 pdfauthor={Mayk Loide Baccaro},
 pdfsubject={Nonreal zeros of finite symmetrized incomplete-gamma approximants to the Riemann Xi function}
}

\newtheorem{theorem}{Theorem}[section]
\newtheorem{proposition}[theorem]{Proposition}
\newtheorem{corollary}[theorem]{Corollary}
\newtheorem{lemma}[theorem]{Lemma}
\theoremstyle{remark}
\newtheorem{remark}[theorem]{Remark}

\newcommand{\C}{\mathbb C}
\newcommand{\R}{\mathbb R}
\newcommand{\XiN}{\Xi_N}
\newcommand{\FNt}{F_{N,t}}
\setlength{\emergencystretch}{2em}

\title[Boundary-constant obstruction]
{A Boundary-Constant Obstruction to Real-Rooted\\
Riemann Incomplete-Gamma Approximants}
\author[Mayk Loide Baccaro]{Mayk Loide Baccaro\\
Independent researcher\\
\href{mailto:m.baccaro7663@student.nu.edu}{\texttt{m.baccaro7663@student.nu.edu}}}
\date{14 August 2026}
\subjclass[2020]{Primary 11M26; Secondary 30D15, 33B20}
\keywords{Riemann Xi function, incomplete gamma function, entire function, real-rootedness, Hadamard factorization}

\begin{document}
\raggedbottom

\begin{abstract}
Let $H_N$ denote the $N$th Fourier partial sum obtained from Riemann's
incomplete-gamma expansion, and let
\[
 c_N=\sum_{n>N}(4\pi n^2-1)e^{-\pi n^2}>0.
\]
For every finite $N\geq1$ and every $0<t\leq1$, we prove that the even entire
function $F_{N,t}=H_N+t c_N$ has a nonreal zero.  The direct arithmetic-index
truncation is the endpoint $F_{N,1}$.  Indeed, $F_{N,t}(x)\to t c_N$ at both
ends of the real axis, while $F_{N,t}(i/2)-t c_N=1/2-c_N\ne0$; moreover,
$F_{N,t}$ has order at most one.  If all zeros were real, the positive
boundary limit would leave only finitely many zeros, and Hadamard
factorization together with evenness would force $F_{N,t}$ to be constant.
P\'olya's theorem already treats the unshifted endpoint $t=0$.  The result
identifies the obstruction created by the positive constant shift and separates it
from Haglund's next-summand zero-motion conjectures.  It has no implication
for the Riemann hypothesis itself.
\end{abstract}

\maketitle

\section{Introduction}

To support a limiting real-rootedness argument, a finite approximating family
needs a corresponding finite-level zero theorem.  Local uniform convergence
alone does not place the approximants' zeros on the real axis.  P\'olya's
Fourier-transform program and Hejhal's refinement show that finite-level zero
behavior depends sensitively on the choice of approximating kernel
\cite{Polya1926,Hejhal1990}.  Haglund later studied incomplete-gamma partial
sums and the geometry of their complex zeros \cite{Haglund2011}.

This note studies direct finite truncations of the symmetrized
incomplete-gamma expansion for the completed zeta function.  These truncations
preserve both functional-equation halves term by term and converge locally
uniformly to Xi.  Conversion to a cosine transform separates an additional
boundary constant.  The exact decomposition derived below is
$\Xi_N=H_N+c_N$.  We place it in the one-parameter family
$F_{N,t}=H_N+t c_N$ and study every positive shift $0<t\leq1$.

The boundary constant forces a nonreal zero at every finite level.

\begin{theorem}\label{thm:main}
For every finite integer $N\geq1$ and every $0<t\leq1$, the function
$\FNt$ defined in~\eqref{eq:shift-family} has at least one nonreal zero.  In
particular, the symmetrized incomplete-gamma approximant
$\XiN=F_{N,1}$ has a nonreal zero for every finite $N$.
\end{theorem}

The argument does not locate a zero numerically.  It combines four global
facts: a positive real-axis boundary limit, nonconstancy, order at most one,
and evenness.  Section~\ref{sec:representation} derives the exact cosine
representation and its endpoint constant.  Section~\ref{sec:boundary}
evaluates the constant by Jacobi's theta transformation.  The growth estimate
and factorization argument appear in Sections~\ref{sec:growth}
and~\ref{sec:zeros}.  Section~\ref{sec:comparison} identifies the unshifted
endpoint covered by P\'olya, distinguishes Haglund's deformation, and states
what the theorem does and does not imply.

\section{The finite incomplete-gamma family}\label{sec:representation}

Write
\[
 \xi(s)=\frac12s(s-1)\pi^{-s/2}\Gamma\!\left(\frac{s}{2}\right)\zeta(s),
 \qquad
 \Xi(z)=\xi\!\left(\frac12+iz\right).
\]
The standard theta-function argument gives a symmetrized incomplete-gamma
expansion for $\xi$; see, for example, \cite[Ch.~1]{Edwards2001} and
\cite{Haglund2011}.  Its direct truncation after the $N$th arithmetic term is
\begin{align}
 \xi_N(s)=\frac12+\frac{s(s-1)}2\sum_{n=1}^N
 \bigg[& (\pi n^2)^{-s/2}
       \Gamma\!\left(\frac{s}{2},\pi n^2\right) \notag\\
 &+(\pi n^2)^{-(1-s)/2}
       \Gamma\!\left(\frac{1-s}{2},\pi n^2\right)\bigg],
 \label{eq:def-xin-s}
\end{align}
where $\Gamma(a,x)=\int_x^\infty t^{a-1}e^{-t}\,dt$.  Define
\begin{equation}\label{eq:def-xin}
 \XiN(z)=\xi_N\!\left(\frac12+iz\right).
\end{equation}
Each $\XiN$ is a real even entire function.  The omitted tail decays
exponentially in $n^2$, uniformly for $z$ in compact sets, so $\XiN\to\Xi$
locally uniformly as $N\to\infty$.

Set
\begin{equation}\label{eq:S-def}
 S_N(u)=\sum_{n=1}^N\exp\!\left(\frac u2-\pi n^2e^{2u}\right),
 \qquad u\geq0,
\end{equation}
and
\[
 I_N(z)=\int_0^\infty S_N(u)\cos(zu)\,du.
\]
Substituting $t=\pi n^2e^{2u}$ in the two incomplete-gamma integrals and
using $s=1/2+iz$ gives
\begin{equation}\label{eq:first-integral}
 \XiN(z)=\frac12-2\left(z^2+\frac14\right)I_N(z).
\end{equation}

Let
\begin{equation}\label{eq:phi-def}
 \Phi_N(u)=S_N''(u)-\frac14S_N(u).
\end{equation}
All derivatives of $S_N$ decay superexponentially at infinity.  Two
integrations by parts therefore yield
\[
 \int_0^\infty \Phi_N(u)\cos(zu)\,du
 =-S_N'(0)-\left(z^2+\frac14\right)I_N(z).
\]
Combining this identity with~\eqref{eq:first-integral} proves the exact
representation
\begin{equation}\label{eq:cosine-representation}
 \XiN(z)=c_N+2\int_0^\infty \Phi_N(u)\cos(zu)\,du,
\end{equation}
where
\begin{equation}\label{eq:c-partial}
 c_N=\frac12+2S_N'(0)
 =\frac12+\sum_{n=1}^N(1-4\pi n^2)e^{-\pi n^2}.
\end{equation}
Put
\begin{equation}\label{eq:shift-family}
 H_N(z)=2\int_0^\infty\Phi_N(u)\cos(zu)\,du,
 \qquad
 F_{N,t}(z)=H_N(z)+t c_N \quad (0\leq t\leq1).
\end{equation}
Then~\eqref{eq:cosine-representation} says exactly that
$\Xi_N=F_{N,1}$.

\section{The positive boundary constant}\label{sec:boundary}

The sign of $c_N$ follows from the theta transformation, not from numerical
cancellation in~\eqref{eq:c-partial}.

\begin{proposition}\label{prop:constant}
For every finite $N\geq1$,
\begin{equation}\label{eq:c-tail}
 c_N=\sum_{n>N}(4\pi n^2-1)e^{-\pi n^2},
 \qquad 0<c_N<\frac12.
\end{equation}
Moreover,
\begin{equation}\label{eq:real-limit}
 \lim_{x\to\pm\infty}\XiN(x)=c_N,
 \qquad x\in\R.
\end{equation}
For $0\leq t\leq1$ one has, more generally,
\begin{equation}\label{eq:shift-limit}
 \lim_{x\to\pm\infty}F_{N,t}(x)=t c_N.
\end{equation}
\end{proposition}

\begin{proof}
Let
\[
 S(u)=\sum_{n=1}^\infty\exp\!\left(\frac u2-\pi n^2e^{2u}\right).
\]
For $\vartheta(x)=\sum_{k\in\mathbb Z}e^{-\pi k^2x}$, Jacobi's identity
$\vartheta(x)=x^{-1/2}\vartheta(1/x)$ gives
\[
 S(u)-S(-u)=-\sinh\!\left(\frac u2\right).
\]
Differentiating at $u=0$ gives $S'(0)=-1/4$, hence
\[
 \frac12+\sum_{n=1}^\infty(1-4\pi n^2)e^{-\pi n^2}=0.
\]
Subtracting~\eqref{eq:c-partial} proves the tail formula
in~\eqref{eq:c-tail}.  Every tail summand is positive.  The strict upper
bound follows directly from~\eqref{eq:c-partial}, since every term in its
finite sum is negative.

Finally, $\Phi_N\in L^1(0,\infty)$.  The Riemann--Lebesgue lemma applied
to~\eqref{eq:cosine-representation} proves~\eqref{eq:real-limit}.
\end{proof}

The shifted functions are not constant.  Indeed, at $z=i/2$ one has $s=0$,
and the factor $s(s-1)$ in~\eqref{eq:def-xin-s} vanishes.  Since
$\Xi_N=H_N+c_N$,
\begin{equation}\label{eq:nonconstant}
 F_{N,t}(i/2)=\frac12-(1-t)c_N,
 \qquad
 F_{N,t}(i/2)-t c_N=\frac12-c_N>0.
\end{equation}

\section{Growth}\label{sec:growth}

\begin{lemma}\label{lem:order}
For each fixed finite $N$ and $t\in[0,1]$, the entire function $F_{N,t}$ has
order at most one.  More precisely, if
$M_{N,t}(r)=\max_{|z|\leq r}|F_{N,t}(z)|$, then
\[
 \log M_{N,t}(r)=O_N\!\bigl(r\log(2+r)\bigr).
\]
\end{lemma}

\begin{proof}
For $|z|\leq r$ and $u\geq0$,
$|\cos(zu)|\leq e^{ru}$.  Since
\[
 0<S_N(u)\leq N\exp\!\left(\frac u2-\pi e^{2u}\right),
\]
the change of variables $t=e^{2u}$ gives
\begin{align*}
 |I_N(z)|
 &\leq \frac N2\int_1^\infty
 t^{r/2-3/4}e^{-\pi t}\,dt\\
 &\leq \frac N2\,\pi^{-(r/2+1/4)}
 \Gamma\!\left(\frac r2+\frac14\right)
\end{align*}
for all sufficiently large $r$.  Stirling's estimate, together with the
quadratic factor in~\eqref{eq:first-integral}, yields the asserted bound.
The identity $F_{N,t}=\Xi_N-(1-t)c_N$ shows that adding the bounded constant
does not change the estimate.  Its order consequence is immediate from the
definition of the order of an entire function.
\end{proof}

\section{Nonreal zeros at every finite level}\label{sec:zeros}

\begin{proof}[Proof of Theorem~\ref{thm:main}]
Fix a finite $N\geq1$ and $0<t\leq1$, and suppose that every zero of
$F_{N,t}$ is real.  Proposition~\ref{prop:constant} and
\eqref{eq:shift-limit} show that $F_{N,t}(x)>0$ for all sufficiently large
$|x|$.  Equation~\eqref{eq:nonconstant} shows that $F_{N,t}$ is not the zero
function, so its zeros are isolated.  Under the supposition, all zeros
therefore lie in a compact real interval and are finite in number.

By Lemma~\ref{lem:order} and Hadamard factorization
\cite[Ch.~2]{Boas1954},
\begin{equation}\label{eq:hadamard}
 F_{N,t}(z)=e^{az+b}P(z)
\end{equation}
for constants $a,b\in\C$ and a polynomial $P$.  Evenness gives
the reflection product
\[
 Q(z)=F_{N,t}(z)F_{N,t}(-z)=e^{2b}P(z)P(-z)
\]
is a polynomial.  If $L=t c_N$, then~\eqref{eq:shift-limit} gives
$Q(x)\to L^2$ as $x\to+\infty$.  A polynomial with a finite limit on the
positive real axis is constant, so $Q\equiv L^2$.  Evenness now gives
$F_{N,t}^2\equiv L^2$.  Since $L>0$, the function $F_{N,t}$ never vanishes;
differentiating shows that $F_{N,t}'\equiv0$.  Hence $F_{N,t}$ is constant,
contrary to~\eqref{eq:nonconstant}.  Thus $F_{N,t}$ has a nonreal zero.
Because $F_{N,t}$ is real and even, its nonreal zeros occur in the
corresponding conjugation and sign symmetry orbit.
\end{proof}

At $t=1$, $F_{N,1}=\Xi_N$ is the direct incomplete-gamma truncation.  Hence
every finite truncation, including $N=1$, has a nonreal zero.  The same
conclusion holds for every positive shift $0<t\leq1$.

\begin{remark}[Local zero motion]\label{rem:motion}
Fix $N$.  If $z_0$ is a simple zero of $F_{N,t_0}$ for some
$0<t_0\leq1$, the holomorphic implicit-function theorem gives a unique local
zero branch $z(t)$ through $(t_0,z_0)$, with
\[
 z'(t_0)=-\frac{c_N}{H_N'(z_0)}.
\]
This identity is local: it gives no global ordering, collision exclusion, or
continuation through a multiple zero.
\end{remark}

\section{Relation to earlier approximants and scope}\label{sec:comparison}

For a single summand
$S_n(u)=\exp(u/2-\pi n^2e^{2u})$, direct differentiation gives
\begin{equation}\label{eq:haglund-normalization}
 2\left(S_n''(u)-\frac14S_n(u)\right)
 =e^{-\pi n^2e^{2u}}
 \left(8\pi^2n^4e^{9u/2}-12\pi n^2e^{5u/2}\right).
\end{equation}
Consequently,~\eqref{eq:shift-family} is exactly the normalization of the
finite Fourier sums in Haglund's equations~(3), (13), and~(14)
\cite{Haglund2011}.  The direct incomplete-gamma truncation differs from
that sum by the positive constant:
\begin{equation}\label{eq:haglund-shift}
 \XiN(z)=H_N(z)+c_N.
\end{equation}

At the unshifted endpoint, P\'olya proved that every finite partial sum of
Riemann's kernel series has infinitely many zeros but only finitely many real
zeros \cite[\S4]{Polya1926}.  Thus $H_N=F_{N,0}$ already has nonreal zeros.
Theorem~\ref{thm:main} covers every positive shift $0<t\leq1$; its boundary
argument requires $t c_N>0$ and therefore does not apply at $t=0$.

Haglund investigated monotonic patterns among complex zeros and reported
numerical zeros without rigorous error bounds \cite{Haglund2011}.  His
numbered conjectures concern the unshifted sums and a next-summand homotopy,
not the additive constant path $H_N+t c_N$.  Hejhal's work concerns other
P\'olya-type kernel modifications \cite{Hejhal1990}.  Neither result addresses
the additive path $H_N+t c_N$, and no monotonicity statement follows from the
local formula in Remark~\ref{rem:motion}.

The conclusion is limited to the functions in~\eqref{eq:shift-family} with
$0<t\leq1$.  It rules out the proof plan that demands real-rootedness for
these finite approximants and then passes to $\Xi$ by local uniform
convergence.  Changing the truncation, adding a nonconstant correction, or
proving a different source-specific zero theorem produces a different
problem.  The existence of nonreal zeros of the approximants has no
implication for whether the zeros of the limiting Riemann Xi function are
real.

\section*{Data and code availability}

The proof is analytic and uses no external data or numerical zero location.
An accompanying Lean~4 formalization checks the finite definitions, theta
identity, boundary limits, nonconstancy witness, growth bound, factorization
contradiction, symmetry orbit, and local-uniform convergence to the classical
completed Xi function, together with the simple-zero local motion stated in
Remark~\ref{rem:motion}.

\section*{Assistance disclosure}

AI systems were used extensively in mathematical derivation, Lean proof
development, literature discovery, organization, and typesetting.  The author
remains responsible for every mathematical statement, proof, citation, and
submission decision.

\bibliographystyle{amsplain}
\bibliography{references}

\end{document}
